\section{R ile çözümlü uygulama}
\label{sec:r-b12}

\textbf{Soru.} Birinci grupta 30 başarılı, 10 başarısız; ikinci grupta
20 başarılı, 20 başarısız gözlem vardır. Pearson bağımsızlık testini,
beklenen frekansları ve Cramér $V$ değerini bulun.

\begin{lstlisting}[style=prismR]
gozlenen <- matrix(c(30, 10, 20, 20), nrow = 2,
                    byrow = TRUE,
                    dimnames = list(
                      grup = c("Birinci", "Ikinci"),
                      sonuc = c("Basarili", "Basarisiz")))
sonuc <- chisq.test(gozlenen, correct = FALSE)
print(sonuc)
sonuc$expected
katkilar <- (gozlenen - sonuc$expected)^2 / sonuc$expected
katkilar
sonuc$residuals
cramer_v <- sqrt(as.numeric(sonuc$statistic) /
                 (sum(gozlenen) * (min(dim(gozlenen)) - 1)))
cramer_v
prop.table(gozlenen, margin = 1)
stopifnot(all(sonuc$expected >= 5))
\end{lstlisting}

\begin{outputguidebox}
\textbf{Çözüm ve kontrol değerleri.} Her satır için beklenen sayılar
$(25,15)$; hücre katkıları $(1;1{,}6667)$'dir. Böylece
$\chi^2(1)=5{,}3333$, $p=0{,}02092$ ve $V=0{,}2582$ olur. Başarı oranları
birinci grupta 0,75, ikinci grupta 0,50'dir. Pearson artıkları ilk satırda
$(1;-1{,}2910)$, ikinci satırda $(-1;1{,}2910)$'dır.
\end{outputguidebox}

\textbf{Yorum.} \texttt{correct = FALSE} bölümdeki düzeltmesiz Pearson
hesabını üretir; ikiye iki tablolarda varsayılan süreklilik düzeltmesiyle aynı
sonuç beklenmez. Bağımsız gözlemler ve uygun örnekleme modeli gerekir;
beklenen frekans denetimi yalnız sayısal koşullardan biridir.

\textbf{Pekiştirme ve yanıt.} Sonuç nesnesinin \texttt{residuals} bileşeni
Pearson artıklarını, \texttt{stdres} bileşeni farklı bir ölçeklendirmeyi verir. Bunları aynı
artık türü gibi raporlamayın. Ki-kare testi ilişkinin nedensel olduğunu göstermez.